\documentclass[../../main.tex]{subfiles}
\begin{document}


{\bf [解]}
$A=(bc-1)(ca-1)(ab-1)$ とおく．以下 $a,b,c\in\mathbb{N}$ に注意する．

\paragraph{(1)}
$A$を展開すると
\begin{align}
A=(abc)^2-(a^2bc+ab^2c+abc^2)+(bc+ca+ab)-1 = (abc)^2-(a+b+c)(abc)+(bc+ca+ab-1)
\end{align}
だから，$A/abc\in\mathbb{Z}$ となる時（$bc+ca+ab-1$）$/abc\in\mathbb{Z}$ つまり $bc+ca+ab-1$ が $abc$ で割り切れる．

\paragraph{(2)}
(1)より $k\in\mathbb{N}$ に対して（$\because bc+ca+ab-1>0$）
\begin{align}
\dfrac{bc+ca+ab-1}{abc}=k \quad \therefore \dfrac{1}{a}+\dfrac{1}{b}+\dfrac{1}{c}-\dfrac{1}{abc}=k \label{eq:1}
\end{align}
とおける．又，$a\in\mathbb{N}$ 及び題意の不等式$1<a<b<c$より
\begin{align}
2\le a<b<c, \quad a,b,c\in\mathbb{N} \label{eq:2}
\end{align}
である．
\cref{eq:2}を用いて\cref{eq:1}の左辺を上から評価すると
\begin{align}
 k = \dfrac{1}{a}+\dfrac{1}{b}+\dfrac{1}{c}-\dfrac{1}{abc}<\dfrac{1}{a}+\dfrac{1}{a}+\dfrac{1}{a}-\dfrac{1}{abc}=\dfrac{1}{a}\left(3-\dfrac{1}{bc}\right)
\end{align}
となり，$k\in\mathbb{N}$ から
\begin{align}
1<\dfrac{1}{a}\left(3-\dfrac{1}{bc}\right) \quad\therefore\ a<3-\dfrac{1}{bc}<3 \quad (\because a,b,c>0)
\end{align}
なる評価を得る．\cref{eq:2}と合わせて 
\begin{align}
 a=2 \label{eq:3}
\end{align}
が必要である．\cref{eq:1}に代入して
\begin{align}
\dfrac{1}{b}+\dfrac{1}{c}-\dfrac{1}{2}\cdot\dfrac{1}{bc}=k-\dfrac{1}{2} \label{eq:4}
\end{align}
を得る．

$a$を得たのと全く同様にして$b$を求めよう．\cref{eq:4}の左辺を上から$b<c$で評価し，右辺を下から$k\ge 1$で評価すると
\begin{align}
 \dfrac{1}{2}\le k-\dfrac{1}{2}<\dfrac{1}{b}\left(2-\dfrac{1}{2c}\right) \quad\therefore\ b<2\left(2-\dfrac{1}{2c}\right)<4 \quad (\because b,c>0)
\end{align}
だから\cref{eq:2}より
\begin{align}
 b = 3 \label{eq:5}
\end{align}
が必要である．\cref{eq:4}に代入すると
\begin{align}
\dfrac{1}{c}-\dfrac{1}{3}\cdot\dfrac{1}{c}=k-\dfrac{1}{2}-\dfrac{1}{3} \quad\therefore\ c=\dfrac{1}{\dfrac{6}{5}k-1} \quad \left(\because k\in\mathbb{N}\text{から}\dfrac{6}{5}k-1>0\right) \label{eq:6}
\end{align}
を得る．ここで\cref{eq:2,eq:5}より$c>3$だから
\begin{align}
 \dfrac{1}{\dfrac{6}{5}k-1}>3 \iff k<\dfrac{25}{18} \quad \left(\because\dfrac{6}{5}k-1>0\right)  
\end{align}
となり，これを満たす自然数$k$は$k=1$のみである．この時\cref{eq:6}より$c=5$である．
以上から
\begin{align}
 (k,a,b,c) = (1,2,3,5)
\end{align}
が必要である．この時実際に$A=630, abc = 30$で確かに$A$は$abc$で割り切れるから題意を満たして十分である．
よって求める$(a,b,c)$の値は
\begin{align}
 (a,b,c) = (2,3,5)
\end{align}
である．

\newpage

\end{document}
